% Generated draft: requires mathematical and source review.
\documentclass[11pt]{article}
\usepackage[margin=25mm]{geometry}
\usepackage{fontspec}
\setmainfont{Latin Modern Roman}
\usepackage{amsmath,amssymb,mathtools,unicode-math}
\setmathfont{Latin Modern Math}
\usepackage{enumitem,needspace,xurl,hyperref}
\setlength{\parindent}{0pt}
\setlength{\parskip}{5pt}
\setlength{\emergencystretch}{3em}
\newcommand{\N}{\mathbb{N}}
\newcommand{\Z}{\mathbb{Z}}
\newcommand{\Q}{\mathbb{Q}}
\newcommand{\R}{\mathbb{R}}
\newcommand{\C}{\mathbb{C}}
\newcommand{\F}{\mathbb{F}}
\newcommand{\A}{\mathbb{A}}
\newcommand{\PP}{\mathbb{P}}
\newcommand{\E}{\mathbb{E}}
\newcommand{\eps}{\varepsilon}
\newcommand{\dd}{\,\mathrm{d}}
\begin{document}
\section*{UnsolvedMath ID: 800786}
\textbf{Source problem code:} catalogue entry
\subsection{Definitions and mathematical statement}
The supplied excerpt suggests a one-complex-variable setting through its discussion of \8501calle--Voronin invariants, but does not explicitly specify the ambient dimension. The following formulation is conditional on that interpretation.

A holomorphic germ at the origin is an equivalence class of holomorphic maps defined on neighborhoods of \(0\in\C\), where two representatives are equivalent if they agree on a neighborhood of \(0\). A germ \(f:(\C,0)\to(\C,0)\) is tangent to the identity if
\[
f(0)=0,\qquad f'(0)=1.
\]
Two such germs \(f\) and \(g\) are holomorphically locally conjugate if there exists a holomorphic germ \(h:(\C,0)\to(\C,0)\) satisfying
\[
h(0)=0,\qquad h'(0)\ne0,\qquad h\circ f=g\circ h.
\]
Here the last equality is equality of germs; the derivative condition ensures that \(h\) has a holomorphic local inverse.

\textbf{Problem OP4 (Abate, 2011).} Give an effective procedure that, for every pair of holomorphic germs \(f,g\) tangent to the identity, decides whether they are holomorphically locally conjugate.

The supplied record does not specify how arbitrary analytic germs are presented to the procedure, what operations on their coefficients are permitted, or what constitutes an effective decision procedure. These choices are required for a precise algorithmic formulation. It also does not specify whether the identity germ itself is excluded.

\textbf{Supplied contextual qualifications.} The excerpt discusses an analogous question for parabolic germs. It reports that the multiplier and a suitably defined parabolic multiplicity form a complete set of topological invariants; adjoining \8501calle's iterative residue gives a complete set of formal invariants; and adjoining suitably adjusted \8501calle--Voronin functional invariants gives a complete set of holomorphic invariants. It describes the parabolic index as another possible quantity, while preferring the iterative residue, and emphasizes the usual practical difficulty of computing the functional invariants. These are contextual assertions from the supplied source, not independently verified assertions here. The excerpt supplies neither the definitions and hypotheses needed for these classifications nor the ensuing generalized problem: it ends with an introduction to that missing question. Thus no precise additional parabolic problem is asserted in this draft.
\subsection{Short English statement}
Abate's Problem OP4 asks for an effective way to decide whether two holomorphic germs tangent to the identity are conjugate by a holomorphic change of coordinates near their fixed point. The formulation above assumes one complex variable, as suggested by the context. The input representation, meaning of effectiveness, and treatment of the identity germ remain unspecified. The registry labels the problem open; that label has not been independently verified.
\subsection{Sources}
\begin{itemize}
\item Local UnsolvedMath registry, permanent ID \(800786\), exact source problem code \texttt{catalogue entry}, title \emph{Local Discrete Holomorphic Dynamics --- Open Problem OP4}. The supplied record attributes the problem to Abate, with proposed year 2011, and labels its status \texttt{open}.
\item Abate, \emph{Open Problems in Local Discrete Holomorphic Dynamics} (2011), Problem OP4. Supplied source link: \url{https://arxiv.org/abs/1107.4876}. The record reports access on 2026-10-16 and extraction method \texttt{arxiv\_latex}; these are supplied provenance statements, not evidence of independent source consultation for this draft.
\end{itemize}
\end{document}
